\section{Mapping and Agreement Records}
\label{supp:mapping}
\subsection{Data Extraction and Mapping}

From the set of included articles, we extracted the data needed to address RQ1. For each proposed methodology, we identified the construction phases and the specific activities associated with each phase. The first author extracted the data and organized them in a spreadsheet. The second author reviewed the extracted data, and when disagreements arose, both researchers discussed the inconsistencies until they reached a consensus.

\added[id=A]{We initially employed an inductive coding approach to analyze and organize the phases and activities extracted from the selected methodologies. The analysis preserved the terminology and process organization reported in the original studies, allowing different activities and process structures to emerge from the literature. During this analysis, we observed substantial convergence between the emerging categories and the phases and activities defined in the EKEBN process proposed by Mendes~\cite{mendes2012using}.}

\added[id=A]{After observing this correspondence, and drawing on the authors' prior familiarity with EKEBN, we adopted its phases and activities as reference categories to facilitate subsequent deductive coding. The coding remained open to additional phases or activities within the checklist's scope. Unmatched activities were recorded under ``Others'' and examined in relation to that scope, rather than treating the reference categories as an exhaustive account of the modeling-project lifecycle. The treatment of these activities is reported with the mapping results below.}

\added[id=A]{The first three authors conducted the deductive mapping of the activities extracted from the methodologies (which we refer to as general activities) to the EKEBN process. This mapping consisted of two stages:}

\begin{itemize}
\item \textit{Mapping general activities to the EKEBN phases.} \added[id=A]{In this stage, we examined whether the activities identified in the analyzed methodologies could be organized according to the EKEBN phases or whether additional phases of the BN construction process emerged from the analysis.} All activities extracted from the methodologies were organized in a spreadsheet, along with the descriptions provided in the articles. Each researcher independently analyzed and classified the activities according to the EKEBN phases: structure building, parameter estimation, and model validation. Activities that did not correspond to any of these phases were assigned to the category ``Others.'' In cases of disagreement, the three researchers discussed the classifications until reaching a consensus. To assess inter-rater agreement, we calculated the Fleiss' Kappa coefficient ($\kappa$)~\cite{fleiss1971measuring}, whose value is presented in Table~\ref{tab:fleiss_kappa_stage1}. According to the interpretation proposed by Landis and Koch~\cite{landis1977measurement}, the researchers achieved an almost perfect level of agreement.

\item \textit{Mapping general activities to the EKEBN activities.} \added[id=A]{In this stage, we examined whether the activities retained within the checklist's scope could be organized according to the activities defined in EKEBN or whether additional BN construction activities emerged from the analysis. The input consisted of the activities grouped by construction phase after the first-stage mapping and scope exclusions described below.} Each researcher analyzed these activities individually and classified them according to the activities defined for each EKEBN phase. Activities related to structure building were classified as ``identification of nodes/variables,'' ``identification of node values/states,'' ``identification of relationships,'' or ``structure evaluation.'' Activities related to parameter estimation were classified according to categories associated with ``CPT construction'' and ``CPT evaluation.'' Finally, activities related to model validation were classified as ``expert-based validation'' or ``data-driven validation.'' \added[id=A]{An ``Others'' category was again available for unmatched activities, but none of the retained activities required it in this second stage.} As in the previous stage, some activities generated classification disagreements. In such cases, the three researchers discussed the divergences until reaching a consensus. The Fleiss' Kappa coefficient ($\kappa$) was also calculated in this stage to assess inter-rater agreement. The values of the coefficient for activity mapping within each phase, along with their interpretation, are presented in Table~\ref{tab:fleiss_kappa_stage2}.
\end{itemize}

\begin{table}[t]
\caption{Fleiss' $\kappa$ coefficient for inter-rater agreement (first stage of the analysis)}
\label{tab:fleiss_kappa_stage1}
\centering
\small
\renewcommand{\arraystretch}{1.15}
\setlength{\tabcolsep}{10pt}

\begin{tabular}{lc}
\toprule
\textbf{$\kappa$ value} & \textbf{Level of agreement} \\
\midrule
0.93 & Almost perfect \\
\bottomrule
\end{tabular}
\end{table}

\begin{table}[t]
\caption{Fleiss' $\kappa$ coefficients for inter-rater agreement by phase (second stage of the analysis)}
\label{tab:fleiss_kappa_stage2}
\centering
\small
\renewcommand{\arraystretch}{1.15}
\setlength{\tabcolsep}{10pt}

\begin{tabular}{lcc}
\toprule
\textbf{Phase} & \textbf{$\kappa$ value} & \textbf{Level of agreement} \\
\midrule
Structure building     & 0.72 & Substantial \\
Parameter estimation  & 1.00 & Almost perfect \\
Model validation      & 0.74 & Substantial \\
\bottomrule
\end{tabular}
\end{table}

\subsection{Mapping Results}

Next, we present the results and analysis of the systematic review. The list of selected articles and the mapping results are available in the supplementary ~\cite{rique2026supplementary}.



To address RQ1, we extracted the phases and activities described in the BN construction process from each methodological article. The extraction followed a descriptive approach that preserved the terminology used by the studies' authors, respecting the original ways in which each paper used a methodology to structure the BN processes. The results revealed substantial heterogeneity in how the methodologies organize the construction process, showing both common phases and activities and elements that capture aspects unique to each proposed approach. Although different methodologies refer to similar activities using distinct terminology (e.g., ``identify variables'' vs. ``identify key factors''), we retained the descriptions exactly as they appear in the original articles. This decision preserves the integrity of each methodology's contribution and enables a more faithful analysis of convergences and divergences in later stages.

The extraction produced a spreadsheet (see~\cite{rique2026supplementary}) containing all phases and activities reported in the analyzed articles. For example, P13 proposed the following phases: ``Identify the objective,'' ``Develop the network structure,'' ``Formalize the model structure,'' ``Quantify the model,'' ``Interrogate the BN model,'' ``Validate the BN model,'' and ``Communicate BN results.'' Within the ``Identify the objective'' phase, the methodology included the activity ``Identify one top-level node.'' Within ``Develop the network structure,'' the proposed activities were ``Identify key factors,'' ``Identify remaining factors,'' ``Identify node states,'' and ``Identify relationships.'' Another example of an article that organized the construction process differently is P17, which defined the following phases: ``Explore problem,'' ``Variables,'' ``Structure,'' ``Parameters,'' ``Explore network,'' and ``Report.'' Within ``Variables,'' the proposed activities were ``Specify BN variables'' and ``Specify discrete states of BN variables.'' Within the ``Structure'' phase, the ``Specify relationships'' activity was included.

The resulting spreadsheet provides an overview of the activities considered across methodologies and serves as the basis for mapping these activities to the EKEBN process. During the mapping process, the need to split certain activities for greater clarity emerged from the researchers' individual analyses and was subsequently discussed and consensually validated during the joint meetings. For example, the activity ``Add, delete, or modify nodes and directed links'' in P16 was separated into ``Add, delete, or modify nodes'' and ``Add, delete, or modify directed links.'' The mapping process aimed to eliminate redundancies, identify gaps, and consolidate a reference process for constructing expert-based or hybrid BN models. 

\added[id=A]{The first deductive mapping stage assigned the general activities extracted from each article to an EKEBN phase or to ``Others.''} Table~\ref{tab:map_stage1_examples} presents examples of this mapping (the complete mapping is documented in the supplementary material~\cite{rique2026supplementary}).

\begin{table}[t]
\caption{Examples of activities mapped to EKEBN phases in the first stage of the analysis}
\label{tab:map_stage1_examples}
\centering
\small
\renewcommand{\arraystretch}{1.2}
\setlength{\tabcolsep}{8pt}

\begin{tabularx}{\linewidth}{>{\raggedright\arraybackslash}X
                              >{\centering\arraybackslash}p{1.2cm}
                              >{\raggedright\arraybackslash}p{2.5cm}}
\toprule
\textbf{Activity} & \textbf{Paper} & \textbf{Phase} \\
\midrule
Determine model objective & P1 & Others \\
Identify the initial set of nodes & P1 & Structure building \\
Expert review of predictions & P2 & Model validation \\
Eliminate cycles & P4 & Structure building \\
Derivation of the conditional probabilities associated with the variables in the map & P10 & Parameter estimation \\
Populate conditional probability tables (CPTs) & P16 & Parameter estimation \\
\bottomrule
\end{tabularx}
\end{table}

\begin{table}[t]
\caption{Activities classified as ``Others'' in the first stage of the analysis}
\label{tab:map_stage1_others}
\centering
\small
\renewcommand{\arraystretch}{1.2}
\setlength{\tabcolsep}{8pt}

\begin{tabularx}{0.70\linewidth}{>{\raggedright\arraybackslash}X
                              >{\centering\arraybackslash}p{1.2cm}}
\toprule
\textbf{Activity} & \textbf{Paper} \\
\midrule
Determine model objective & P1 \\
Perform decision analysis & P1 \\
Validate objective & P13 \\
Create the modes of results communication & P13 \\
Identify the hypotheses suggested by the problem/questions & P17 \\
Identify the items of evidence most pertinent to those hypotheses & P17 \\
Develop report & P17 \\
Determine quantification approach & P25 \\
Collate relevant information & P25 \\
Determine assessment baseline & P25 \\
Perform validation of the objective & P25 \\
Identifying the stakeholders & P29 \\
Defining the question the BN is supposed to answer & P29 \\
Identifying suitable data and information to build the model & P29 \\
Selecting experts & P29 \\
Selecting appropriate tools & P29 \\
\bottomrule
\end{tabularx}
\end{table}

\added[id=A]{The first-stage mapping identified activities that did not correspond directly to the EKEBN phases. These activities were recorded as ``Others'' (Table~\ref{tab:map_stage1_others}) and were not retained in the checklist. The checklist's scope concerns activities that directly specify, parameterize, or evaluate the BN itself, assuming that a BN has been selected as the modeling approach. It does not cover the entire modeling-project lifecycle. The exclusions concern the following types of activities:}

\begin{itemize}
\item \added[id=A]{Project-level preparation, context definition, and approach selection. These include defining or checking project objectives, identifying stakeholders and available information sources, and selecting experts, tools, or an overall modeling approach. Such activities support the modeling project but are outside the checklist's focus on constructing and validating the BN.}

\item \added[id=A]{Model use and communication. These include establishing a baseline for decision analyses, using the model to support decisions, communicating results, and preparing reports. They concern the application and communication of the model rather than its construction or validation and were not included as checklist activities.}
\end{itemize}

\added[id=A]{This is a functional scope boundary, not a requirement that activities occur in a fixed temporal order or be exclusive to BNs. Identifying variables that become BN nodes and validating the BN remain within scope even though analogous activities occur in other modeling approaches. Conversely, excluding project-level preparation or model-use activities does not imply that they are unnecessary for a successful modeling project.}

\added[id=A]{P25~\cite{chakraborty2015multifaceted} illustrates the distinction. We interpreted ``Determine quantification approach'' as selecting and motivating the overall quantitative modeling approach; the corresponding passage discusses alternatives such as regression and factor analysis before introducing BNs. It was classified as ``Others.'' The separate activity ``Quantify model,'' which describes estimating BN probability distributions, was retained under parameter estimation and subsequently mapped to CPT construction. Thus, the classification distinguishes approach selection from parameter estimation, even though P25 places both activities within its Quantification stage.}

\added[id=A]{After these scope exclusions, the retained activities were organized into the three EKEBN phases. These phase-level groups provided the input for the second mapping stage.}

\added[id=A]{In the second mapping stage, all retained activities could be accommodated by the EKEBN activities, and none required the ``Others'' category. Table~\ref{tab:map_stage2_examples} presents examples per phase (the complete mapping is available in the supplementary material~\cite{rique2026supplementary}). This finding supports the use of EKEBN as an organizing framework for the activities retained in this synthesis. It does not independently establish EKEBN's completeness, because the finding is conditional on the reviewed corpus and the checklist's scope.}

\begin{table}[t]
\caption{Examples of activities mapped to EKEBN activities in the second stage of the analysis}
\label{tab:map_stage2_examples}
\centering
\small
\renewcommand{\arraystretch}{1.2}
\setlength{\tabcolsep}{8pt}

\begin{tabularx}{\linewidth}{>{\raggedright\arraybackslash}X
                              >{\centering\arraybackslash}p{1.2cm}
                              >{\raggedright\arraybackslash}p{3.2cm}}
\toprule
\textbf{Activity} & \textbf{Paper} & \textbf{EKEBN Activity} \\
\midrule

\multicolumn{3}{l}{\textbf{Structure building}} \\
Introduce synthetic variables & P1 & Identification of nodes/variables \\
Manage mutual exclusivity & P1 & Identification of node values/states \\
Build initial causal graph & P4 & Identification of relationships \\
Revise causal graph with experts & P4 & Structure evaluation \\

\addlinespace[6pt]
\multicolumn{3}{l}{\textbf{Parameter estimation}} \\
Learning parameters with insufficient data & P7 & CPT construction \\
Populate conditional probability tables (CPTs) & P16 & CPT construction \\
Evaluate CPTs & P16 & CPT evaluation \\

\addlinespace[6pt]
\multicolumn{3}{l}{\textbf{Model validation}} \\
Performance evaluation (in terms of discrimination, calibration, and accuracy) & P7 & Data-driven validation \\
Validate results with experts & P13 & Expert-based validation \\

\bottomrule
\end{tabularx}
\end{table}

\begin{tcolorbox}[
    colback=gray!5,
    colframe=black,
    boxrule=0.5pt,
    arc=2pt,
    left=6pt,
    right=6pt,
    top=6pt,
    bottom=6pt
]
\textbf{RQ1 --- What are the fundamental activities involved in constructing expert-based or hybrid BNs?}

The analysis revealed significant heterogeneity in how methodologies describe and group BN construction activities. Some approaches focus narrowly on activities directly related to network construction, while others also include more general modeling steps, such as defining objectives, scope, and context. \added[id=A]{Within the scope retained for the checklist, the activities were organized into:}
(i) structure building (including identification of variables, states, relationships, and structure evaluation),
(ii) parameter estimation (including CPT construction and evaluation), and
(iii) model validation (expert-based or data-driven).

\added[id=A]{EKEBN accommodated the retained activities across the analyzed methodologies and provided the organizing framework for the checklist. This coverage concerns the construction and validation activities included in the synthesis, rather than the entire modeling-project lifecycle.}
\end{tcolorbox}
